% ============================================================================
% sec_classification.tex — V8: Alternative FL classification methods
% Addresses external referee Concern 2: ad hoc classification rule
% ============================================================================

\subsection{Classification Robustness: IQR Rule vs.\ Model-Based Alternatives} % V8
\label{sec:classification_robustness}

A concern with the IQR threshold rule is that it may be ad hoc and
could capture inefficient firms rather than cover bidders. We compare
the IQR rule with two model-based alternatives on the same CADE
ground truth.

\paragraph{Structural classifier.}
Using the estimated Regime~2 parameters ($\hat{\epsilon} = 0.83$,
$\hat{\sigma}_c = 1.19$), we compute a log-likelihood ratio score for
each always-loser firm: the likelihood of the firm's mean bid spread
under the cover-bid distribution relative to the genuine-loser spread
distribution. Firms with high scores are more likely cover bidders
according to the structural model. This score achieves
AUC~$= 0.56$---barely above random---because the bid-spread
distributions of genuine losers and cover bidders overlap
substantially. The structural model identifies cover bidding at the
\emph{distributional} level (Regime~2 vs.\ Regime~1) but its
parameters are insufficiently discriminative at the \emph{individual
firm} level.

\paragraph{Machine-learning classifier.}
We train a random forest (500 trees, 5-fold cross-validation to
prevent overfitting) on seven firm-level features: number of
participations, number of tenders, mean/SD/median bid spread,
fraction of bids above the winner, and bid-level CV. The
cross-validated AUC is $0.90$---strong, but \emph{below} the IQR
rule's $0.94$ (DeLong $p = 0.04$).

\paragraph{Interpretation.}
Table~\ref{tab:classification_comparison} and
Figure~\ref{fig:classification_comparison} summarize the comparison.
The IQR rule dominates both the structural classifier and the random
forest. This result has two implications. First, the threshold rule
is not merely convenient; it captures the most discriminative feature
of cover bidding---\emph{abnormally high participation despite zero
wins}---which ML cannot improve upon by adding bid-level features.
Second, the structural model's bid-level parameters are complementary
(they identify the \emph{mechanism}) but not substitutes (they do not
improve \emph{detection}) for the participation-based screen.

\begin{table}[htbp]
\centering
\caption{FL Classification: IQR Rule vs.\ Model-Based Alternatives}
\label{tab:classification_comparison}
\begin{threeparttable}
\small
\begin{tabular}{lccc}
\toprule
Classifier & AUC & Data required & DeLong $p$ \\
\midrule
IQR rule (this paper) & 0.937 & Participation count & --- \\
Structural score & 0.562 & Bid values & 0.0000 \\
Random Forest (7 features) & 0.904 & Bid values + firm chars & 0.0416 \\
Combined (IQR + structural) & 0.894 & Both & \\
Random classifier & 0.500 & None & --- \\
\bottomrule
\end{tabular}
\begin{tablenotes}
\small
\item \textit{Notes:} AUC computed on always-loser firms with valid
bid-level features. Ground truth: co-participation with CADE-convicted
cartelists. IQR rule: participation count (higher $=$ more suspicious).
Structural score: log-likelihood ratio of mean bid spread under Regime~2
cover-bid distribution ($\hat{\epsilon} = 0.83$, $\hat{\sigma}_c = 1.19$)
vs.\ genuine-loser spread distribution.
Random Forest: 500 trees trained on 7 firm-level features (participations,
mean/SD/median spread, \% above winner, bid CV).
DeLong $p$-value: null that AUC difference from IQR rule equals zero.
\end{tablenotes}
\end{threeparttable}
\end{table}


\begin{figure}[htbp]
\centering
\includegraphics[width=0.78\textwidth]{images/fig_classification_comparison.pdf}
\caption{ROC curves: IQR rule vs.\ model-based classifiers.
The IQR rule (AUC~$= 0.94$) dominates the random forest with
7 bid-level features (AUC~$= 0.90$; DeLong $p = 0.04$) and the
structural likelihood-ratio score (AUC~$= 0.56$). Ground truth:
co-participation with CADE-convicted cartelists.}
\label{fig:classification_comparison}
\end{figure}
